% GENERATED, edit the French unit or the translation table
% GENERATED from content/data/enrichment.json by mxbuild.enrichment; do not edit.
\mxunit{enrichment-bisection-solutions}{Bracketing a root by bisection}
\mxsubtitle{Upper high school · analysis and algorithms \textperiodcentered{} Worked solutions \textperiodcentered{} 12}

\begin{note}
\textbf{Given information.} For exercise 1: seek a zero of $f(x)=x^{2}-2$ on $\intff{1}{2}$. After two halvings, the retained interval is $\intff{\num{1,25}}{\num{1,5}}$.
\end{note}

\begin{exercise}
Carry out a third step of the example.
\end{exercise}
\begin{solution}
The midpoint is $\num{1,375}$, whose square is $\num{1,890625}<2$. The new interval is $\intff{\num{1,375}}{\num{1,5}}$.
\end{solution}

\begin{exercise}
Starting with width 1, how many halvings guarantee width at most $\frac{1}{1000}$?
\end{exercise}
\begin{solution}
$2^{9}=512<1000\leq 1024=2^{10}$, so 10 halvings suffice and 9 do not.
\end{solution}

\begin{exercise}
Why does the sign change of $f(x)=\frac{1}{x}$ between $-1$ and 1 not prove the existence of a zero?
\end{exercise}
\begin{solution}
The function is undefined at 0 and is not continuous on $\intff{-1}{1}$. The theorem’s hypothesis fails; $\frac{1}{x}$ is never zero.
\end{solution}

\begin{exercise}
A root lies in $\intff{\num{1,4}}{\num{1,5}}$. What maximum error is guaranteed by midpoint $\num{1,45}$?
\end{exercise}
\begin{solution}
The half-width is $\frac{\num{1,5}-\num{1,4}}{2}=\num{0,05}$.
\end{solution}

\begin{problem}[Investigate, justify, conclude]
For $f(x)=x^{3}-x-1$ on $\intff{1}{2}$, prove existence and uniqueness of a root $\alpha $. Carry out three bisection steps. How many halvings guarantee that the final midpoint approximates $\alpha $ with error at most $\num{0,001}$?
\end{problem}
\begin{solution}
$f(1)=-1$ and $f(2)=5$, and f is continuous. Also $f'(x)=3x^{2}-1\geq 2$ on $\intff{1}{2}$, so f is strictly increasing. Midpoints $\num{1,5}$, $\num{1,25}$, $\num{1,375}$ give $\intff{1}{\num{1,5}}$, $\intff{\num{1,25}}{\num{1,5}}$, $\intff{\num{1,25}}{\num{1,375}}$ because f takes values $\num{0,875}$, $-\num{0,296875}$ and $\num{0,224609375}$. After n halvings, midpoint error is at most $\frac{1}{2^{n+1}}$. Thus $2^{n+1}\geq 1000$, requiring $n\geq 9$.
\end{solution}

